\section{Technical Problems and Solutions}
\label{sec:technical_problems_and_solutions}

During the development, architectural refactoring, and feature expansion of the 2D Platformer engine, several critical technical challenges emerged across the \textbf{Animation Subsystem}, \textbf{Finite State Machine (FSM)}, \textbf{Object-Oriented Hierarchy}, \textbf{World Physics Engine}, and \textbf{State Persistence Architecture}. This section details the most prominent engineering bottlenecks encountered and the mathematical, algorithmic, and architectural solutions designed to resolve them.

% =========================================================================
% SECTION 4.1
% =========================================================================
\subsection{Sprite Origin Misalignment and Center-of-Mass Jittering}
\label{subsec:sprite_misalignment}

\subsubsection{Problem Formulation}
When integrating high-resolution, non-uniform sprite sheets for advanced character forms (e.g., \textit{Epic Archmage Flash}, \textit{Fire Mario}, and \textit{Kitsune Luigi}), significant visual anomalies were observed during gameplay:
\begin{enumerate}
    \item \textbf{Horizontal Jerking / Positional Sliding:} Transitioning between discrete animation states (e.g., $\text{Idle} \rightarrow \text{Run} \rightarrow \text{Special Skill}$) caused characters to abruptly shift horizontally across the terrain.
    \item \textbf{Limb and Prop Truncation:} Attempting to eliminate jitter by tightly cropping sprite bounding boxes inadvertently severed character extremities (such as Mario's outstretched casting arm or Luigi's bushy fox tail).
\end{enumerate}

The root cause originated from the generic origin assignment implemented in the rendering utility:
\begin{equation}
    \mathbf{O}_{\text{generic}} = \left( \frac{W_{\text{rect}}}{2}, \; H_{\text{rect}} \right)
\end{equation}
Because character silhouettes within tight rectangular bounds ($W_{\text{rect}}$) are geometrically asymmetric, the bounding box midpoint $\frac{W_{\text{rect}}}{2}$ deviated substantially from the character's physical vertical anchor axis (the physical center of mass, $CM_x$), generating an instantaneous visual displacement upon every state transition.

\subsubsection{Engineering Solution}
Rather than compromising asset integrity through tight cropping or polluting the core engine with hardcoded per-frame offsets, an automated \textbf{Mathematical Center-of-Mass ($CM$) Padding Algorithm} was formulated and executed via image-processing scripts.

\paragraph{1. Center of Mass Computation:} For an arbitrary frame sub-region of width $W$ and height $H$, the horizontal center of mass is computed across all opaque body pixels:
\begin{equation}
    CM_x = \frac{\sum_{y=Y_{\min}}^{Y_{\max}} \sum_{x=X_{\min}}^{X_{\max}} x \cdot \mathbb{I}(\alpha(x, y) > \tau)}{\sum_{y=Y_{\min}}^{Y_{\max}} \sum_{x=X_{\min}}^{X_{\max}} \mathbb{I}(\alpha(x, y) > \tau)}
\end{equation}
where $\alpha(x, y) \in [0, 255]$ represents the pixel alpha channel, $\tau = 25$ is the opacity threshold, and $\mathbb{I}(\cdot)$ is the indicator function.

\paragraph{2. Symmetric Geometric Padding:} Given the calculated anchor $CM_x$ and the true visual bounds $[X_{\min}, X_{\max}]$, symmetric transparent space is injected to force the geometric midpoint to coincide with $CM_x$:
\begin{equation}
    X_{\text{rect}}, W_{\text{rect}} = 
    \begin{cases} 
        \left( X_{\min}, \; 2 \cdot (CM_x - X_{\min}) \right) & \text{if } CM_x \ge \frac{X_{\min} + X_{\max}}{2} \\
        \left( 2 \cdot CM_x - X_{\max}, \; 2 \cdot (X_{\max} - CM_x) \right) & \text{if } CM_x < \frac{X_{\min} + X_{\max}}{2}
    \end{cases}
\end{equation}
This formulation guarantees that the origin $\frac{W_{\text{rect}}}{2}$ remains strictly aligned with the character's physical foot position ($100\%$ jitter elimination) while preserving all extended visual effects and staff/tail animations.

% =========================================================================
% SECTION 4.2
% =========================================================================
\subsection{Finite State Machine (FSM) Oscillation and Frame Freezing}
\label{subsec:fsm_oscillation}

\subsubsection{Problem Formulation}
During locomotion testing, characters exhibited a severe defect where holding the movement keys caused the sprite to freeze on a single static frame while translating across the floor (appearing to ``skate'' without playing the stride animation).

Code inspection of the locomotion state machine (\texttt{RunState.cpp}) revealed an unguarded velocity threshold check:
\begin{equation}
    \text{if } (|v_x| < v_{\text{stop}}) \implies \text{Transition to } \mathtt{IdleState}
\end{equation}
When accelerating from rest, initial horizontal velocity $v_x$ initiates at $0.0\,\text{px/s}$, which is strictly below $v_{\text{stop}} = 8.0\,\text{px/s}$. Consequently, on every frame:
\begin{enumerate}
    \item \texttt{IdleState} detected keypress $\rightarrow$ transitioned to \texttt{RunState}.
    \item On the immediate subsequent tick, \texttt{RunState} detected $|v_x| < 8.0\,\text{px/s}$ $\rightarrow$ prematurely reverted to \texttt{IdleState}.
    \item This high-frequency state oscillation continuously triggered \texttt{Hero::setState()}, resetting the animation accumulator ($\mathtt{currentFrameIndex} = 0$), permanently locking the character on Frame 0.
\end{enumerate}

\subsubsection{Engineering Solution}
The transition guard was redesigned to decouple player intentionality from instantaneous kinematic magnitude by constructing a composite boolean invariant:
\begin{equation}
    \text{Transition}(\mathtt{RunState} \rightarrow \mathtt{IdleState}) \iff \left( \neg \mathtt{pressLeft} \land \neg \mathtt{pressRight} \right) \land \left( |v_x| < v_{\text{stop}} \right)
\end{equation}
This condition ensures that as long as directional input remains active, the FSM remains locked in \texttt{RunState}, allowing uninterrupted animation frame delta accumulation and fluid multi-frame locomotion.

% =========================================================================
% SECTION 4.3
% =========================================================================
\subsection{Elimination of Hardcoded Ground Constraints via Dynamic Gravity Kinematics}
\label{subsec:dynamic_gravity}

\subsubsection{Problem Formulation}
In early engine iterations, physical entities (including power-up items such as \textit{Mushrooms} and \textit{Flowers}, dynamic throwables such as \textit{Potions}, and spawned \textit{Enemies}) relied on static, hardcoded vertical offsets ($Y$-coordinates) mapped to flat ground planes. This introduced severe structural fragility:
\begin{itemize}
    \item When deploying multi-tier levels (such as underground rooms, floating platforms, and elevated pipe structures in World 1-1 and World 1-3), entities spawned with fixed parameters either hovered in mid-air or intersected solid terrain blocks.
    \item Kinematic item trajectories (e.g., mushrooms popping out of Question Blocks) lacked realistic arc landing mechanics across irregular terrain contours.
\end{itemize}

\subsubsection{Engineering Solution}
To enforce map-agnostic behavior, static spatial assumptions were eliminated in favor of a \textbf{Unified Dynamic Gravity and AABB Collision Resolution Subsystem} within \texttt{WorldPhysicsSystem}:

\paragraph{1. Kinematic Velocity Integration:} Every dynamic entity updates vertical velocity under constant gravitational acceleration:
\begin{align}
    v_y(t + \Delta t) &= \min\left(v_y(t) + g \cdot \Delta t, \; v_{\text{terminal}}\right) \\
    y(t + \Delta t) &= y(t) + v_y(t + \Delta t) \cdot \Delta t
\end{align}
where $g = 1568\,\text{px/s}^2$ and $v_{\text{terminal}} = 960\,\text{px/s}$ prevents collision tunneling through thin tile boundaries.

\paragraph{2. Continuous Ground Snapping:} During vertical collision passes against the terrain grid, ground contact is resolved dynamically:
\begin{equation}
    \text{if } \text{CollisionSide} == \mathtt{Top} \implies 
    \begin{cases}
        y_{\text{resolved}} = y_{\text{tile\_top}} - H_{\text{hitbox}} \\
        v_y = 0 \\
        \mathtt{isGrounded} = \text{true}
    \end{cases}
\end{equation}
This mathematical resolution guarantees that all characters, dropped items, and projectiles dynamically conform to arbitrary terrain topography with zero manual configuration.

% =========================================================================
% SECTION 4.4
% =========================================================================
\subsection{Polymorphic Multi-Tier Hero Evolution and Texture Routing}
\label{subsec:polymorphic_evolution}

\subsubsection{Problem Formulation}
Procedural texture loading in \texttt{HeroForm} subclasses previously invoked \texttt{hero->loadTexture(baseTexturePath)} directly upon entry. When integrating heroes with independent asset sets (e.g., \textit{Flash} utilizing \texttt{thunderflash2.png} or \textit{Luigi} utilizing \texttt{KitsuneLuigi.png}), procedural transitions caused texture coordinate mismatches and skipped growth forms, violating the \textbf{Open/Closed Principle (OCP)}.

\subsubsection{Engineering Solution}
The evolution subsystem was refactored using the \textbf{State/Strategy Pattern}:
\begin{enumerate}
    \item \textbf{Polymorphic Texture Dispatch:} Overrode \texttt{Hero::loadTexture()} in specialized classes to dynamically resolve resource bindings based on active form state.
    \item \textbf{Decoupled Multi-Tier Scaling:} Implemented independent rendering scale metrics:
    \begin{equation}
        S_{\text{render}} = 
        \begin{cases} 
            S_{\text{small}} & \text{if } \mathtt{Form} == \text{Small} \quad (32\times 32\,\text{px hitbox}) \\
            S_{\text{giant}} & \text{if } \mathtt{Form} == \text{Giant} \quad (32\times 64\,\text{px hitbox}) \\
            S_{\text{special}} & \text{if } \mathtt{Form} == \text{Fire/Special} \quad (32\times 64\,\text{px hitbox})
        \end{cases}
    \end{equation}
    \item \textbf{Uniform 3-Tier Hierarchy:} Standardized progression across all playable heroes: $\text{Small Form} \xrightarrow{\text{Mushroom}} \text{Giant Form} \xrightarrow{\text{Flower}} \text{Tier-3 Special Form}$.
\end{enumerate}

% =========================================================================
% SECTION 4.5
% =========================================================================
\subsection{World State Serialization and Dynamic Deserialization (Save/Load)}
\label{subsec:save_load_architecture}

\subsubsection{Problem Formulation}
Implementing persistent world serialization (\texttt{F5} Quick Save, \texttt{F9} Quick Load) required capturing a heterogeneous runtime state containing:
\begin{itemize}
    \item Dynamic entities (enemies with remaining state timers, active projectiles, roaming items).
    \item Spatial world modifications (permanently hit question blocks, shattered brick blocks removed from the collision matrix).
    \item Player inventory, form type, invincibility I-frames, and HUD progress metrics.
\end{itemize}

\subsubsection{Engineering Solution}
A \textbf{Base Level Reload $\rightarrow$ State Overlay Pipeline} was engineered within \texttt{SaveManager} using JSON serialization:
\begin{enumerate}
    \item \textbf{Phase 1: Deterministic Base Reload:} \texttt{LevelRuntime::reload()} re-instantiates the raw TMX/TMJ tile matrix and baseline colliders.
    \item \textbf{Phase 2: Dynamic State Overlay:}
    \begin{itemize}
        \item \textbf{Tile Matrix Mutation:} Querying coordinates $(T_x, T_y)$ against serialized \texttt{destroyedBlocks} updates the interactive vertex mesh via \texttt{setTileID("Interactive", tx, ty, 0)}.
        \item \textbf{Polymorphic Entity Reconstruction:} \texttt{EnemyFactory} and \texttt{EnemyStateFactory} reconstruct surviving enemies with precise velocity vectors and remaining AI timers (e.g., maintaining a Koopa shell's retreat countdown).
        \item \textbf{Hero State Restoration:} Re-injects exact health, coin count, form name, and remaining invulnerability duration.
    \end{itemize}
\end{enumerate}

\subsection{Dynamic Spatial Partitioning via Playable Regions and Context-Aware Theming}
\label{subsec:playable_regions}
\subsubsection{Problem Formulation}
In traditional monolithic 2D platformers, levels often encompass multiple distinct sub-environments within a single continuous tilemap (e.g., the sunny Overworld surface alongside subterranean coin chambers or deep castle arenas). Managing such multi-room levels through a single global coordinate system created three critical architectural bottlenecks:
\begin{enumerate}
    \item \textbf{Camera Viewport Bleeding:} A naive player-following camera easily drifted out of room boundaries, exposing unrendered void areas, black borders, or adjacent hidden rooms on the screen.
    \item \textbf{Static Pit-Death Invalidation:} In early iterations, the ``pit-fall'' death boundary was defined as a single hardcoded global horizontal line ($Y_{\text{death}} = Y_{\text{global\_bottom}}$). When a player entered an elevated room or a deep underground chamber situated far below the main level surface, this static threshold caused either \textit{instantaneous false deaths upon entering lower rooms} or \textit{unbounded falling through infinite space} without triggering character death.
    \item \textbf{Heterogeneous Environmental Theming:} Objects, terrain tiles, and background visual styles (e.g., vibrant Overworld vs. dark blue Underground vs. molten Castle) required dynamic context-aware rendering based strictly on the player's active containment zone.
\end{enumerate}
\subsubsection{Engineering Solution}
To resolve these spatial and visual conflicts, the engine implemented a **Data-Driven Dynamic Spatial Partitioning Architecture** based on decoupled \texttt{PlayableRegion} bounding zones and active \texttt{Theme} identifiers parsed dynamically from Tiled map triggers:
\paragraph{1. Spatial Zone Partitioning:} Each distinct chamber within the map is encapsulated by a geometric bounding volume:
\begin{equation}
    \mathcal{R}_k = \left\{ (x, y) \in \mathbb{R}^2 \;\middle|\; X_k^{\min} \le x \le X_k^{\max}, \; Y_k^{\min} \le y \le Y_k^{\max} \right\}
\end{equation}
The engine tracks the hero's position $\mathbf{P}_{\text{hero}}(t)$ and dynamically resolves the active region:
\begin{equation}
    \mathcal{R}_{\text{active}}(t) = \left\{ \mathcal{R}_k \;\middle|\; \mathbf{P}_{\text{hero}}(t) \in \mathcal{R}_k \right\}
\end{equation}
\paragraph{2. Constrained Camera Viewport Clamping:} The game camera dynamically clamps its center $(C_x, C_y)$ to ensure that the rendering viewport of dimension $(W_{\text{cam}}, H_{\text{cam}})$ never exposes out-of-bounds void:
\begin{align}
    C_x &= \text{clamp}\left( P_{\text{hero}, x}, \; X_k^{\min} + \frac{W_{\text{cam}}}{2}, \; X_k^{\max} - \frac{W_{\text{cam}}}{2} \right) \\
    C_y &= \text{clamp}\left( P_{\text{hero}, y}, \; Y_k^{\min} + \frac{H_{\text{cam}}}{2}, \; Y_k^{\max} - \frac{H_{\text{cam}}}{2} \right)
\end{align}
\paragraph{3. Dynamic Context-Aware Kill Planes:} Rather than relying on a static coordinate, the lethal pit threshold ($Y_{\text{kill}}$) for both the Hero and Enemies is computed dynamically relative to the active region's bottom boundary:
\begin{equation}
    Y_{\text{kill}}(\mathcal{R}_k) = Y_k^{\max} + \delta_{\text{threshold}}
\end{equation}
where $\delta_{\text{threshold}}$ represents the vertical grace offset before triggering \texttt{Hero::die()} or enemy entity deallocation. When traveling through pipes into an underground cavern, $Y_{\text{kill}}$ instantly re-binds to the subterranean floor baseline.
\paragraph{4. Context-Driven Visual and Audio Theming:} Each $\mathcal{R}_k$ is mapped to a dedicated environmental descriptor (\texttt{Theme::Overworld}, \texttt{Theme::Underground}, \texttt{Theme::Castle}):
\begin{equation}
    \mathcal{T}_{\text{render}} = f(\mathcal{R}_{\text{active}}) \implies 
    \begin{cases}
        \text{Background Asset \& Clear Color Binding} \\
        \text{Interactive Object / Palette Shader Dispatch} \\
        \text{Background Music (BGM) Stream Switching}
    \end{cases}
\end{equation}
\subsubsection{Resulting Impact}
This architecture completely eliminates camera visual bleeding across room transitions, guarantees mathematical accuracy for pit-fall death detection across arbitrary multi-tier map layouts, and enables rich, heterogeneous environmental theming within a single seamless level file.   